\chapter{La plateforme : architecture et observatoire}
\label{ch:plateforme}

\questionchapitre{Un entrepôt ne se consulte pas. Quelle architecture permet d'exploiter la
donnée consolidée sans enfermer les méthodes développées dans l'application qui les héberge ?}

\section{Une contrainte d'architecture posée d'emblée}

La plateforme comporte trois couches : une interface web, un service applicatif, et une
bibliothèque métier en Python. La contrainte structurante porte sur la troisième :
\textbf{la bibliothèque métier ne comporte aucune dépendance à l'interface}.

Cette règle a une conséquence pratique déterminante pour un travail de recherche. Le code de
préparation des données, de calibration des modèles métier, d'entraînement, d'explicabilité et
de génération de contrefactuels s'exécute indifféremment depuis le service web, depuis un
script ou depuis un carnet de calcul. Les méthodes développées ne sont donc pas prisonnières de
l'application : si l'interface était abandonnée, elles resteraient utilisables. C'est la
condition pour que les développements survivent au projet, exigence explicite d'une mission dont
la finalité est d'outiller des travaux ultérieurs.

\begin{figure}[H]
\centering
\begin{tikzpicture}[
  font=\sffamily,
  bande/.style={draw=black!70, fill=white, minimum width=10.6cm, minimum height=0.85cm,
                align=center, font=\sffamily\footnotesize},
  coeur/.style={draw=black, very thick, fill=black!3, minimum width=10.6cm,
                minimum height=1.9cm, align=center, font=\sffamily\footnotesize},
  mod/.style={draw=black!50, fill=white, minimum width=2.35cm, minimum height=0.7cm,
              align=center, font=\sffamily\tiny},
  res/.style={draw=black!55, fill=white, minimum width=3.3cm, minimum height=0.8cm,
              align=center, font=\sffamily\scriptsize},
  fl/.style={-{Stealth[length=2.2mm]}, draw=black!70, thick},
  note/.style={font=\sffamily\scriptsize\itshape, color=black!60},
]

\node[bande] (spa) {\textbf{Interface web} \\[1pt]
                    \scriptsize cartes, graphiques, suivi d'entraînement};
\node[bande, below=0.75cm of spa] (api) {\textbf{Service applicatif} \\[1pt]
                    \scriptsize API et flux d'événements};
\node[coeur, below=0.75cm of api] (core) {};
\node[anchor=north, font=\sffamily\footnotesize\bfseries, yshift=-3pt] at (core.north)
     {Bibliothèque métier};
\node[anchor=north, font=\sffamily\scriptsize, yshift=-17pt] at (core.north)
     {aucune dépendance d'interface};

\node[mod] (m1) at ($(core.center)+(-3.75,-0.42)$) {Préparation\\des données};
\node[mod] (m2) at ($(core.center)+(-1.25,-0.42)$) {Modèles\\métier};
\node[mod] (m3) at ($(core.center)+(1.25,-0.42)$)  {Prévision\\profonde};
\node[mod] (m4) at ($(core.center)+(3.75,-0.42)$)  {Explicabilité\\contrefactuels};

\node[res, below=0.95cm of core, xshift=-3.65cm] (cache) {Cache de réponses};
\node[res, below=0.95cm of core, draw=black, very thick] (entrepot)
     {Entrepôt \\[1pt] \scriptsize lecture SQL};
\node[res, below=0.95cm of core, xshift=3.65cm] (registre) {Registre d'expériences};

\draw[fl] (spa) -- (api);
\draw[fl] (api) -- (core);
\draw[fl] (core.south) -- (entrepot.north);
\draw[fl] (core.south west) ++(1.1,0) -- ++(0,-0.3) -| (cache.north);
\draw[fl] (core.south east) ++(-1.1,0) -- ++(0,-0.3) -| (registre.north);

\node[note, below=0.14cm of entrepot] {dépôt de l'entrepôt, chapitre~\ref{ch:entrepot}};

\end{tikzpicture}
\caption{Architecture de la plateforme. La bibliothèque métier est appelable depuis le service
web, un script ou un carnet de calcul, sans modification. C'est cette propriété qui rend les
développements réutilisables au-delà de l'application.}
\label{fig:archi-plateforme}
\end{figure}

\section{L'Observatoire}

\subsection{Le besoin}

Un entrepôt bien construit reste inerte tant qu'il n'est pas regardé. Deux publics le
demandaient, pour des raisons différentes.

Les hydrogéologues ont besoin de situer une station dans son contexte : quelle nappe
observe-t-elle, comment se comporte-t-elle par rapport à ses voisines, où en est-elle dans son
historique, et que faisait le climat pendant ce temps. Les modélisateurs ont besoin de choisir
leurs cas d'étude en connaissance de cause, car une chronique trop courte, trop lacunaire ou
visiblement perturbée par un prélèvement ne fera pas un bon cas d'apprentissage. Repérer ces
situations avant d'entraîner un modèle évite d'attribuer à la méthode ce qui relève de la
donnée.

\subsection{Organisation}

L'interface se divise en deux volets. Le premier couvre les \textbf{nappes et les rivières} :
exploration cartographique des stations, superposition des entités aquifères et d'un découpage
en secteurs hydrogéologiques, indicateurs de situation portés par les stations, comparaison
entre stations avec conservation de la sélection.

Le second couvre le \textbf{climat} selon trois vues : une situation cartographique nationale,
une vue par maille donnant l'historique depuis 1950 avec la pluie comparée à la normale, les
indices sur plusieurs fenêtres, la table des épisodes de sécheresse et un export tabulaire,
enfin une comparaison interannuelle. Les deux volets communiquent : depuis la carte des nappes,
un lien contextuel ouvre le climat de la maille correspondante.

\subsection{Une règle de conception}

Les couches cartographiques obéissent à une règle explicite : \emph{soit un indicateur
réellement standardisé, soit une valeur réelle affichée comme telle, jamais un intermédiaire
construit pour l'occasion}.

Concrètement, les indices standardisés sont cartographiés sur l'échelle à sept classes et les
grandeurs journalières brutes sur une échelle absolue. En revanche, les cumuls mensuels absolus
de pluie, de température ou d'évapotranspiration ne sont pas proposés comme couches : ils
apparaissent comme chiffres exacts dans le panneau de détail d'une maille.

La raison est méthodologique. Cartographier une grandeur absolue sur une palette suppose un
seuillage, donc une norme implicite, et cette norme n'étant pas une référence statistique
établie, elle serait une invention de l'outil. L'utilisateur lirait une carte de sécheresse là
où il n'y a qu'un choix de bornes. Refuser cette couche coûte une fonctionnalité apparente et
évite une erreur d'interprétation systématique.

\subsection{Trois choix techniques}

\paragraph{L'Observatoire ne recalcule rien.} Ses points d'accès sont des lectures directes des
tables analytiques. Aucune statistique n'est calculée à la volée. La conséquence est double :
les temps de réponse sont ceux d'une base correctement indexée, et surtout la valeur affichée
est exactement celle qui a été validée au chapitre~\ref{ch:indicateurs}, sans second calcul
susceptible de diverger.

\paragraph{Les réponses sont mises en cache.} Les vues nationales portent sur des volumes
importants et changent peu ; leur mise en cache pour vingt-quatre heures réduit la charge sans
affecter la fraîcheur perçue. Le corollaire opérationnel est documenté : tout déploiement qui
modifie la forme ou la valeur des réponses doit purger le cache, sous peine de servir des
valeurs périmées pendant une journée.

\paragraph{Le point d'entrée est unique.} Le service qui sert l'interface assure aussi le relais
vers l'API, les en-têtes de sécurité et la limitation de débit. Une architecture antérieure
comportait un second relais, redondant, qui n'était jamais démarré en production, de sorte que
les protections qu'il portait n'étaient en pratique jamais appliquées. Sa suppression au profit
d'une configuration unique a supprimé cet écart entre l'intention et le déploiement réel.

\begin{remarque}[title={Une protection non déployée est une protection absente}]
Le cas mérite d'être retenu au-delà de son détail technique : un composant configuré mais jamais
démarré fournit la croyance qu'une garantie existe, sans la garantie. La consolidation en un
point d'entrée unique vaut moins pour sa simplicité que pour la disparition de cet écart.
\end{remarque}

\section{Accès, sécurité et données personnelles}

La plateforme héberge des comptes nominatifs et les travaux de leurs titulaires, ce qui impose
des exigences distinctes de celles de l'entrepôt.

L'accès est provisionné par un administrateur : il n'existe ni inscription libre ni
authentification déléguée. Les mots de passe sont hachés par un algorithme à coût mémoire
paramétrable et ne sont détenus en clair par personne, administrateur compris. La session
repose sur un jeton en cookie non accessible au script, révocable côté serveur par
désactivation du compte ou incrément d'un compteur de version relu à chaque requête. Deux rôles
existent, l'administration et l'usage courant, ce dernier ne donnant accès qu'aux objets dont
l'utilisateur est propriétaire.

Un garde-fou mérite mention : la suppression du dernier administrateur actif est refusée. En
l'absence d'inscription libre, elle rendrait l'application définitivement inaccessible par
l'interface. La contrainte est portée par le code plutôt que par une consigne.

Enfin, les journaux d'authentification sont conservés trois cent soixante-cinq jours puis
supprimés automatiquement, et l'utilisateur peut effacer son compte et les données associées
depuis son espace personnel.

\section{Environnements et déploiement}

Deux environnements complètement séparés coexistent sur le serveur de calcul, l'un suivant la
branche principale, l'autre la branche de développement, avec des volumes, des bases, des
ports et des artefacts distincts. Un modèle entraîné ou un compte créé dans l'un n'affecte
jamais l'autre. Cette séparation est ce qui autorise à expérimenter sur une chaîne en évolution
sans compromettre l'instance servant aux utilisateurs.

La production est répartie : l'interface est déployée sur l'infrastructure mutualisée de
l'établissement, le calcul et les bases restent sur le serveur équipé d'un processeur
graphique. Une seule exigence réseau conditionne l'ensemble, l'ouverture de la route entre les
deux.

\begin{acquisouvert}
\textbf{Acquis.} Une architecture à trois couches dont le cœur métier est indépendant de
l'interface, un observatoire alimenté directement par les indicateurs validés, une gestion des
accès et des données personnelles conforme aux attentes d'un service hébergeant des travaux
nominatifs, et une séparation stricte des environnements.

\textbf{Ouvert.} L'Observatoire dépend entièrement de l'entrepôt et ne se dégrade pas
gracieusement en son absence : il répond en erreur plutôt que d'afficher un état partiel. Ce
comportement est assumé pour un outil scientifique et serait à revoir pour une diffusion large.
\end{acquisouvert}
